Mapa indexu ekonomické závislosti starých osob (OLDDEP) v~zemích EU27 dokumentuje, že ČR se v~horizontu 2050 řadí mezi státy se středně nadprůměrným demografickým stárnutím --- přičemž projekce do roku 2070 předpokládají přiblížení k~hodnotám Německa nebo Itálie.
Demografický tlak na průběžný důchodový systém je přitom citlivý nejen na absolutní počet seniorů, ale i~na výši průměrné mzdy ze které se odvádí pojistné: systém, jehož příjmy stagnují kvůli nízké mzdové dynamice, je dvojnásob ohrožen rostoucím dependency ratio.
Pracovní produktivita a~mzdový růst tedy nejsou jen ekonomickým, ale i~fiskálním imperatívem --- každé procentní zvýšení reálné mzdy přeložené do vyšší pojistné základny prodlužuje solventnost systému více než administrativní parametrické změny.
Kolektivní smlouvy s~pravidelnou indexací mezd jsou v~tomto kontextu institucí s~pozitivní externalitou pro celý sociální stát: zvyšují pojistnou základnu bez nutnosti zvyšovat sazby --- a~tím zpomalují tempo, kterým rostoucí dependency ratio nahlodává důchodový systém (srov.~obr.~\ref{fig:cz_pension_sp_ratio}).

Mapa indexu ekonomické závislosti starých osob (OLDDEP) v~zemích EU27 dokumentuje, že ČR se v~horizontu 2050 řadí mezi státy se středně nadprůměrným demografickým stárnutím --- přičemž projekce do roku 2070 předpokládají přiblížení k~hodnotám Německa nebo Itálie.
Demografický tlak na průběžný důchodový systém je přitom citlivý nejen na absolutní počet seniorů, ale i~na výši průměrné mzdy ze které se odvádí pojistné: systém, jehož příjmy stagnují kvůli nízké mzdové dynamice, je dvojnásob ohrožen rostoucím dependency ratio.
Pracovní produktivita a~mzdový růst tedy nejsou jen ekonomickým, ale i~fiskálním imperatívem --- každé procentní zvýšení reálné mzdy přeložené do vyšší pojistné základny prodlužuje solventnost systému více než administrativní parametrické změny.
Kolektivní smlouvy s~pravidelnou indexací mezd jsou v~tomto kontextu institucí s~pozitivní externalitou pro celý sociální stát: zvyšují pojistnou základnu bez nutnosti zvyšovat sazby --- a~tím zpomalují tempo, kterým rostoucí dependency ratio nahlodává důchodový systém (srov.~obr.~\ref{fig:cz_pension_sp_ratio}).

Mapa indexu ekonomické závislosti starých osob (OLDDEP) v~zemích EU27 dokumentuje, že ČR se v~horizontu 2050 řadí mezi státy se středně nadprůměrným demografickým stárnutím --- přičemž projekce do roku 2070 předpokládají přiblížení k~hodnotám Německa nebo Itálie.
Demografický tlak na průběžný důchodový systém je přitom citlivý nejen na absolutní počet seniorů, ale i~na výši průměrné mzdy ze které se odvádí pojistné: systém, jehož příjmy stagnují kvůli nízké mzdové dynamice, je dvojnásob ohrožen rostoucím dependency ratio.
Pracovní produktivita a~mzdový růst tedy nejsou jen ekonomickým, ale i~fiskálním imperatívem --- každé procentní zvýšení reálné mzdy přeložené do vyšší pojistné základny prodlužuje solventnost systému více než administrativní parametrické změny.
Kolektivní smlouvy s~pravidelnou indexací mezd jsou v~tomto kontextu institucí s~pozitivní externalitou pro celý sociální stát: zvyšují pojistnou základnu bez nutnosti zvyšovat sazby --- a~tím zpomalují tempo, kterým rostoucí dependency ratio nahlodává důchodový systém (srov.~obr.~\ref{fig:cz_pension_sp_ratio}).

Mapa indexu ekonomické závislosti starých osob (OLDDEP) v~zemích EU27 dokumentuje, že ČR se v~horizontu 2050 řadí mezi státy se středně nadprůměrným demografickým stárnutím --- přičemž projekce do roku 2070 předpokládají přiblížení k~hodnotám Německa nebo Itálie.
Demografický tlak na průběžný důchodový systém je přitom citlivý nejen na absolutní počet seniorů, ale i~na výši průměrné mzdy ze které se odvádí pojistné: systém, jehož příjmy stagnují kvůli nízké mzdové dynamice, je dvojnásob ohrožen rostoucím dependency ratio.
Pracovní produktivita a~mzdový růst tedy nejsou jen ekonomickým, ale i~fiskálním imperatívem --- každé procentní zvýšení reálné mzdy přeložené do vyšší pojistné základny prodlužuje solventnost systému více než administrativní parametrické změny.
Kolektivní smlouvy s~pravidelnou indexací mezd jsou v~tomto kontextu institucí s~pozitivní externalitou pro celý sociální stát: zvyšují pojistnou základnu bez nutnosti zvyšovat sazby --- a~tím zpomalují tempo, kterým rostoucí dependency ratio nahlodává důchodový systém (srov.~obr.~\ref{fig:cz_pension_sp_ratio}).

\input{texparts/python/old_age_dependency_map}



